% Section 4, System.

% ---- full-width pipeline figure (the paper's one two-column float) ---------
\begin{figure*}[!t]
\centering
% ===========================================================================
% fig-pipeline.tex — the \sysname pipeline, end to end.
%
% MOUNTS AS figure*  (FULL TEXT WIDTH, spanning BOTH columns).  This is the
% one deliberate exception to the single-column house rule: the diagram is a
% seven-stage left-to-right narrative and is unreadable at column width.
% Wrap it as \begin{figure*}[!t] ... \end{figure*}; wrapper/caption/label
% live in the section file.  BODY ONLY (tikzpicture in a \resizebox).
% Requires % ===========================================================================
% fig-style.tex — shared colour / marker scheme for EVERY main-paper figure.
% \input this ONCE from the preamble of main.tex (after tikz + pgfplots).
%
% House rule: the palette is defined here and nowhere else. Each series also
% gets a distinct marker AND line style so the figures survive grayscale
% printing.
%
%   \sysname / deployed .. cSys   (blue!70!black)    mark *
%   opus / claude ........ cOpus  (red!70!black)     mark square*
%   gpt .................. cGpt   (teal!70!black)    mark triangle*
%   kimi ................. cKimi  (violet!70!black)  mark diamond*
%   flash ................ cFlash (orange!70!black)  mark pentagon*
%   greedy / baseline .... cBase  (gray)             mark o
% ===========================================================================
% The pipeline diagram needs the calc library for coordinate arithmetic; the
% rest of the libraries are already loaded in main.tex's preamble.
\usetikzlibrary{calc}
% Hatching, so bar charts stay readable in grayscale.
\usetikzlibrary{patterns}

\colorlet{cSys}{blue!70!black}
\colorlet{cOpus}{red!70!black}
\colorlet{cGpt}{teal!70!black}
\colorlet{cKimi}{violet!70!black}
\colorlet{cFlash}{orange!70!black}
\colorlet{cBase}{gray}

% Common axis furniture: footnotesize ticks/legend, small titles, framed
% legend in gray!50, light horizontal grid.
\pgfplotsset{
  paperaxis/.style={
    width=0.80\linewidth,
    height=5.2cm,
    tick label style={font=\footnotesize},
    label style={font=\footnotesize},
    title style={font=\small},
    legend style={font=\footnotesize, draw=gray!50, fill=white,
                  fill opacity=0.9, text opacity=1, rounded corners=1pt},
    every axis plot/.append style={line width=0.7pt},
    grid style={gray!25, line width=0.3pt},
    axis line style={gray!60},
    tick style={gray!60},
  },
  papermark/.style={mark size=2.6pt},
}
 in the preamble.
%
% PROVENANCE
%   Schematic only — this figure carries NO measured numbers, so there is
%   nothing to verify against a report. Stage names follow the system as
%   described in internal records: verify-then-strip, the search-episode and
%   handoff split, and the deployed router as the blend of task text and
%   \modelname hidden states.
%
%   Layout note: drawn at ~19cm and \resizebox'd to the \linewidth of a
%   figure* (~18.2cm in IEEEtran conference), i.e. scaled by ~0.96. The
%   smallest type in the figure is \scriptsize (7pt), so the smallest labels
%   land at ~6.7pt on paper; the \footnotesize body rows land at ~7.7pt.
%
%   Colour contract (fig-style.tex): cSys tints mark every component we
%   build or train (searcher, gate, router); neutral grays mark the data
%   artefacts that flow between them (task, handoff, patch, score); the four
%   fixer colours (cGpt/cOpus/cFlash/cKimi) are the ones those models carry
%   in every other figure. Shape, fill weight and stroke weight repeat those
%   distinctions, so the figure survives grayscale printing.
% ===========================================================================
\resizebox{\linewidth}{!}{%
\begin{tikzpicture}[
  x=1cm, y=1cm,
  % ---- panels ----------------------------------------------------------
  panel/.style={rounded corners=2.4pt, line width=0.7pt, anchor=center},
  sysbox/.style={panel, draw=cSys!55, fill=cSys!6},
  artbox/.style={panel, draw=gray!55, fill=gray!4},
  pill/.style={draw=gray!55, fill=gray!8, rounded corners=5pt,
               line width=0.7pt, align=center, inner sep=2pt},
  % ---- type ------------------------------------------------------------
  kick/.style={font=\scriptsize\scshape, text=gray!45!black, anchor=south},
  systitle/.style={font=\small\bfseries, text=cSys!80!black,
                   anchor=north west, inner sep=0pt, align=left},
  arttitle/.style={font=\small\bfseries, text=gray!15!black,
                   anchor=north west, inner sep=0pt, align=left},
  body/.style={font=\footnotesize, text=gray!35!black,
               anchor=north west, inner sep=0pt, align=left},
  micro/.style={font=\scriptsize, text=gray!40!black,
                anchor=north west, inner sep=0pt, align=left},
  note/.style={font=\footnotesize\itshape, text=gray!45!black},
  % ---- flow ------------------------------------------------------------
  flow/.style={-{Stealth[length=2.0mm,width=1.5mm]}, gray!62,
               line width=0.9pt, rounded corners=1.5pt},
  fallback/.style={-{Stealth[length=2.0mm,width=1.5mm]}, gray!58,
                   line width=0.75pt, rounded corners=2pt,
                   dash pattern=on 1.5mm off 1.1mm},
  hair/.style={line width=0.5pt},
]

% =========================================================================
% Geometry.  All stage panels are 3.0cm tall and centred on y=0; the band
% runs left to right with a uniform 2.8mm gap for the flow arrows.  Panel
% widths and centres (total drawing width 19.04cm):
%   task   1.50 @ 0.75  | search 2.70 @  3.13 | handoff 2.80 @  6.16
%   gate   2.80 @ 9.24  | route  2.70 @ 12.27 | pool    3.06 @ 15.43
%   tail   1.80 @ 18.14
% =========================================================================
\def\Ph{3.0}

% ---- stage kickers -------------------------------------------------------
\node[kick] at (3.13,1.64) {Search};
\node[kick] at (6.16,1.64) {Handoff};
\node[kick] at (9.24,1.64) {Gate};
\node[kick] at (12.27,1.64) {Route};
\node[kick] at (15.43,1.64) {Fix};

% =========================================================================
% 1. TASK
% =========================================================================
\node[pill, minimum width=1.5cm, minimum height=1.6cm] (task) at (0.75,0)
  {{\small\bfseries Task}\\[3pt]{\footnotesize\color{gray!35!black} issue text}};

% =========================================================================
% 2. SEARCH EPISODE
% =========================================================================
\node[sysbox, minimum width=2.7cm, minimum height=\Ph cm] (search) at (3.13,0) {};
\node[systitle, text width=2.4cm] (st)
  at ([shift={(2.6mm,-2.8mm)}]search.north west) {\modelname};
\coordinate (sa) at ([yshift=-1.6mm]st.south west);
\draw[hair, cSys!30] (sa) -- ([xshift=-2.6mm]search.east |- sa);
\node[body, text width=2.4cm] at ([yshift=-2.4mm]sa)
  {explores the repo\\[1.5pt] sampled decoding\\[1.5pt] bounded turns};

% search-trace motif: a short walk that ends on a committed file
\begin{scope}[shift={(2.20,-1.16)}]
  \draw[cSys!45, line width=0.6pt]
    (0,0.30) -- (0.62,0.30) -- (0.62,0) -- (1.42,0) -- (1.42,0.30) -- (1.86,0.30);
  \fill[cSys!25] (0,0.30)    circle (1.5pt);
  \fill[cSys!45] (0.62,0)    circle (1.5pt);
  \fill[cSys!70] (1.42,0.30) circle (1.5pt);
  \fill[cSys]    (1.86,0.30) circle (2.1pt);
\end{scope}

% =========================================================================
% 3. STRUCTURED HANDOFF  (document glyph, three field rows)
% =========================================================================
\node[artbox, minimum width=2.8cm, minimum height=\Ph cm] (hand) at (6.16,0) {};
\node[arttitle, text width=2.5cm] (ht)
  at ([shift={(2.6mm,-2.8mm)}]hand.north west) {Structured handoff};
\coordinate (ha) at ([yshift=-1.6mm]ht.south west);
\draw[hair, gray!35] (ha) -- ([xshift=-2.6mm]hand.east |- ha);

\begin{scope}[shift={(5.01,0.30)}]
  \def\dw{2.30} \def\dh{1.58} \def\df{0.28}
  \fill[white] (0,0) -- (\dw-\df,0) -- (\dw,-\df) -- (\dw,-\dh) -- (0,-\dh) -- cycle;
  \draw[gray!60, line width=0.6pt]
    (0,0) -- (\dw-\df,0) -- (\dw,-\df) -- (\dw,-\dh) -- (0,-\dh) -- cycle;
  \draw[gray!60, line width=0.6pt] (\dw-\df,0) -- (\dw-\df,-\df) -- (\dw,-\df);
  \foreach \i/\lab in {1/implicated files, 2/repo notes, 3/repro test}{
    \fill[cSys!45, rounded corners=0.4pt]
      (0.18,-0.06-0.44*\i+0.05) rectangle ++(0.115,0.115);
    \node[micro, text=gray!25!black] at (0.40,-0.06-0.44*\i+0.145) {\lab};
  }
\end{scope}

% =========================================================================
% 4. VERIFY-THEN-STRIP SANDBOX GATE  (funnel)
% =========================================================================
\node[sysbox, minimum width=2.8cm, minimum height=\Ph cm] (gate) at (9.24,0) {};
\node[systitle, text width=2.4cm] (gt)
  at ([shift={(2.6mm,-2.8mm)}]gate.north west) {Verify-then-strip};
\coordinate (ga) at ([yshift=-1.6mm]gt.south west);
\draw[hair, cSys!30] (ga) -- ([xshift=-2.6mm]gate.east |- ga);
\node[micro] at ([yshift=-2.2mm]ga) {replays the repro claim};

\begin{scope}[shift={(9.24,-0.30)}]
  \fill[cSys!12] (-0.60,0.34) -- (0.60,0.34) -- (0.17,-0.10) -- (0.17,-0.40)
              -- (-0.17,-0.40) -- (-0.17,-0.10) -- cycle;
  \draw[cSys!65, line width=0.7pt]
    (-0.60,0.34) -- (0.60,0.34) -- (0.17,-0.10) -- (0.17,-0.40)
    -- (-0.17,-0.40) -- (-0.17,-0.10) -- cycle;
  \node[micro, text=cSys!70!black] at (-1.14,-0.56) {$\checkmark$~replays: kept};
  \node[micro, text=gray!50!black] at (-1.14,-0.88) {$\times$~fails: stripped};
\end{scope}

% =========================================================================
% 5. RESUME ROUTER
% =========================================================================
\node[sysbox, minimum width=2.7cm, minimum height=\Ph cm] (route) at (12.27,0) {};
\node[systitle, text width=2.4cm] (rt)
  at ([shift={(2.6mm,-2.8mm)}]route.north west) {R\'esum\'e router};
\coordinate (ra) at ([yshift=-1.6mm]rt.south west);
\draw[hair, cSys!30] (ra) -- ([xshift=-2.6mm]route.east |- ra);
\node[body, text width=2.4cm] at ([yshift=-2.4mm]ra)
  {task text $+$ \modelname\ hidden states};

% the résumé shelf it compares against
\begin{scope}[shift={(11.07,-1.14)}]
  \draw[gray!40, fill=white, line width=0.5pt, rounded corners=0.6pt]
    (0.10,0.10) rectangle ++(0.40,0.50);
  \draw[gray!52, fill=white, line width=0.5pt, rounded corners=0.6pt]
    (0.05,0.05) rectangle ++(0.40,0.50);
  \draw[cSys!60, fill=cSys!8, line width=0.6pt, rounded corners=0.6pt]
    (0,0) rectangle ++(0.40,0.50);
  \foreach \i in {1,2,3}{
    \draw[cSys!45, line width=0.4pt] (0.08,0.50-0.125*\i) -- (0.32,0.50-0.125*\i);}
  \node[micro, anchor=west, text width=1.65cm] at (0.62,0.28)
    {vs.\ one r\'esum\'e per fixer};
\end{scope}

% =========================================================================
% 6. FIXER POOL
% =========================================================================
\node[artbox, minimum width=3.06cm, minimum height=\Ph cm] (pool) at (15.43,0) {};
\node[arttitle, text width=2.5cm] (pt)
  at ([shift={(2.6mm,-2.8mm)}]pool.north west) {Fixer pool};
\coordinate (pa) at ([yshift=-1.6mm]pt.south west);
\draw[hair, gray!35] (pa) -- ([xshift=-2.6mm]pool.east |- pa);

% chips: {row}{colour}{name}{0 = not chosen, 1 = chosen, 2 = open slot}
\foreach \i/\col/\nm/\sel in {%
    1/cGpt/\fxgpt/0, 2/cOpus/\fxopus/1, 3/cFlash/\fxflash/0, 4/cKimi/\fxkimi/0,
    5/gray/{+ new fixer}/2}{
  \coordinate (c\i) at (14.10,0.24-0.40*\i+0.40);
  \ifcase\sel
    % 0 — in the pool, not selected this time
    \draw[draw=gray!35, fill=white, line width=0.55pt, rounded corners=1.6pt]
      (c\i) rectangle ++(2.66,0.34);
    \fill[\col!55] ([shift={(0.06,0.05)}]c\i) rectangle ++(0.09,0.24);
    \node[anchor=west, font=\footnotesize, text=gray!42!black,
          inner sep=0pt] at ([shift={(0.28,0.17)}]c\i) {\nm};
  \or
    % 1 — the fixer this task is routed to
    \draw[draw=\col!70, fill=\col!10, line width=0.85pt, rounded corners=1.6pt]
      (c\i) rectangle ++(2.66,0.34);
    \fill[\col] ([shift={(0.06,0.05)}]c\i) rectangle ++(0.09,0.24);
    \node[anchor=west, font=\footnotesize\bfseries, text=\col!72!black,
          inner sep=0pt] at ([shift={(0.28,0.17)}]c\i) {\nm};
  \or
    % 2 — the open slot: a r\'esum\'e is the whole onboarding cost
    \draw[draw=gray!45, dash pattern=on 0.9mm off 0.8mm, line width=0.55pt,
          rounded corners=1.6pt] (c\i) rectangle ++(2.66,0.34);
    \node[anchor=west, font=\scriptsize\itshape, text=gray!52!black,
          inner sep=0pt] at ([shift={(0.16,0.17)}]c\i) {\nm: r\'esum\'e only};
  \fi
}
\node[anchor=east, font=\scriptsize, text=cOpus!72!black, inner sep=0pt]
  at ([shift={(-0.03,0.17)}]c2) {$\blacktriangleright$};

% =========================================================================
% 7. PATCH -> OFFICIAL SCORING
% =========================================================================
\node[pill, minimum width=1.8cm, minimum height=0.8cm] (patch) at (18.14,0)
  {{\small\bfseries Patch}};
\node[pill, minimum width=1.8cm, minimum height=1.0cm] (score) at (18.14,-1.32)
  {{\small\bfseries Official}\\[2pt]{\small\bfseries scoring}};

% =========================================================================
% Flow
% =========================================================================
\draw[flow] (task)   -- (search);
\draw[flow] (search) -- (hand);
\draw[flow] (hand)   -- (gate);
\draw[flow] (gate)   -- (route);
\draw[flow] (route)  -- (pool);
\draw[flow] (pool)   -- (patch);
\draw[flow] (patch)  -- (score);

\node[note, anchor=north, align=center, text width=3.3cm]
  at (15.43,-1.72) {the chosen fixer gets the task $+$ surviving handoff};

% =========================================================================
% Fallback: no handoff, the router still fires
% =========================================================================
\draw[fallback] (search.south) -- (3.13,-2.30) -- (12.27,-2.30) -- (route.south);
\node[note, anchor=south] at (7.70,-2.24)
  {no handoff: the chosen fixer proceeds from the issue text alone};

\end{tikzpicture}}

\caption{\textbf{The \sysname pipeline.}
\modelname explores the repository and emits a structured handoff, which a
sandbox gate verifies before a r\'esum\'e router selects one of four frontier
fixers.  The dashed path is the fallback: when no handoff is produced, the
chosen fixer proceeds from the issue text alone.  Adding a new fixer requires
only a r\'esum\'e, not retraining.}
\label{fig:pipeline}
\end{figure*}

\section{The \sysname System}
\label{sec:system}

Figure~\ref{fig:pipeline} traces a single task through the pipeline.  The task
enters a search phase, where \modelname explores the repository and writes a
structured handoff.  A sandbox gate then verifies the handoff's reproduction
claims before a r\'esum\'e router selects one fixer from a pool of frontier
models.  That fixer receives the task and the surviving handoff, produces a
patch, and submits it to official scoring.  The entire trajectory is a single
pass: no parallel sampling, no cross-model escalation mid-task.  When the
searcher produces no handoff at all, the chosen fixer proceeds from the issue
text alone; this fallback defines the floor as fixer-solo performance.

\subsection{The Searcher and Its Handoff}
\label{sec:searcher}

\modelname is a 7B model built on Qwen2.5-Coder~\citep{qwen25coder}, trained
exclusively for the search phase (\S\ref{sec:training}).  Given a task, it
explores the repository, localizes the implicated files, attempts to write a
failing reproduction test, and then produces a handoff document before stopping.
The handoff is a structured artifact kept short by design, typically
about a page of text (4\,KB).  It contains implicated
files with line regions, ranked by confidence; a reproduction attempt specifying
a file, command, and observed output; dead ends the searcher already tried; and
free-form repository notes.  A verified example appears in
Appendix~\ref{apx:handoff-example}.

When the search exhausts its turn budget without committing to a handoff, a
single extra generation step demands one.  Such handoffs are tagged
\emph{forced} and are never pooled with spontaneous ones anywhere in this paper.

\subsection{Verify-then-Strip}
\label{sec:verify-strip}

The searcher's reproduction claims are not trusted.  Before any fixer sees a
handoff, a sandbox replays the claimed reproduction command against the
unpatched repository.  A claim that does not genuinely fail is stripped: both
the test file and the claim itself are removed from the handoff.  A verified
claim, by contrast, is materialized so the fixer can use it directly.
Calibration revealed that most reproduction claims are in fact false
(\S\ref{sec:calibration}), and \S\ref{sec:results} quantifies the guard's
effect at benchmark scale.  Injection is blanket: every routed fixer receives
the surviving handoff identically.

\subsection{The R\'esum\'e Router}
\label{sec:router}

\begin{table}[!tbp]
\centering
\caption{\textbf{The r\'esum\'e router, end to end.}  Each fixer's
r\'esum\'e stores two embedding centroids and a base rate, computed from
public per-task outcomes.  A logistic head per feature space scores
$P(\text{solve})$; the router walks the pool in cheap-first order,
stopping at the first fixer clearing~$\theta$.  Adding a new fixer
requires only 25--50 public outcomes and no retraining.}
\label{tab:router-recipe}
% ===========================================================================
% tab-router-recipe.tex — the deployed resume router, end to end (§4).
% BODY ONLY (tabular). Wrapper/caption/label live in the section file.
% Provenance: code/router/router_spec.json (embedder pin, resume definition,
% feature recipe, hidden-state key, head, route rule and cheap order).
%
% Notes for editing:
%   - The embedder is frozen infrastructure: one exact model and revision
%     computes every embedding in the system and is never fine-tuned.
%   - Features are listed uncentered and are computed per feature space.
%   - Wording rule: phrase onboarding as "adding a new FIXER requires no
%     retraining", never a bare "without retraining".
% ===========================================================================
\begin{tabular}{@{}ll@{}}
\toprule
Component & Specification \\
\midrule
Embedder      & Qwen3-Embedding-0.6B \\
              & frozen \\
              & 4096-token cap, $L_2$-normalized \\
\addlinespace
R\'esum\'e     & per fixer: mean embedding of \\
              & solved, mean of failed, base rate \\
              & from 25--50 public outcomes \\
\addlinespace
Features      & $\cos(x,s)$, $\cos(x,f)$, \\
              & $\cos(x,s){-}\cos(x,f)$, $p$, \\
              & $\cos(x,s){\cdot}p$, \\
              & $(\cos(x,s){-}\cos(x,f)){\cdot}p$ \\
              & uncentered, per feature space \\
\addlinespace
Hidden state  & \modelname\ layer $-4$, pre-decode \\
              & final-position state (3584-d) \\
\addlinespace
Scorer        & logistic regression per space, \\
              & uniform probability average, \\
              & $\theta{=}0.30$, cheapest-adequate \\
\addlinespace
New fixer     & 25--50 public outcomes; \\
              & no retraining of \modelname, \\
              & the embedder, or the scorer \\
\bottomrule
\end{tabular}

\end{table}

Each fixer in the pool is summarized by a r\'esum\'e built from 25--50 public
per-task outcomes.  A r\'esum\'e stores three quantities: the mean embedding of
tasks the fixer solved, the mean embedding of tasks it failed, and a base solve
rate.  Table~\ref{tab:router-recipe} lists the full specification.

At routing time the task is embedded in two complementary feature spaces.  Its
text is encoded by a frozen off-the-shelf embedder~\citep{qwen3embedding}.  Its
content as \modelname experienced it is captured through the searcher's
3{,}584-dimensional pre-decode hidden state, drawn from the fourth-from-last
layer at the final token position.  In each feature space the router computes
uncentered cosine similarities between the task and every r\'esum\'e's solved
and failed centroids, their difference, and the base rate; a single logistic
regression per feature space, shared across fixers, then scores
$P(\text{solve})$ for each fixer from its r\'esum\'e-relative features.  The
routing head blends the
task-text and hidden-state feature sets by uniformly averaging their predicted
probabilities.

Routing walks the pool in cheap-first order and assigns the task to the first
fixer whose predicted probability clears a caution threshold~$\theta$; if none
clears it, the task falls to a designated anchor model (\S\ref{sec:setup}).  Logistic regression was chosen
over a multi-layer perceptron during calibration, where the simpler scorer won
repeatedly (\S\ref{sec:calibration}).

\FloatBarrier

\subsection{Extensibility}
\label{sec:extensibility}

Adding a new fixer requires no retraining.  A new model's r\'esum\'e consists
of two averaged embeddings and a base rate, computed from its public outcomes.
The embedder is frozen and \modelname is untouched; \S\ref{sec:results}
presents a measured discussion.
